% ============================================================
\chapter*{Pr\'eface}
\addcontentsline{toc}{chapter}{Pr\'eface}
% ============================================================

\section*{Pourquoi 90\,\% des projets ML \'echouent en production}

Un constat r\'ecurrent dans l'industrie est alarmant~: selon Gartner (2022),
seulement 54\,\% des mod\`eles de machine learning (ML) passent du prototype
\`a la production, et parmi ceux-ci, beaucoup sont retir\'es dans les premiers
mois. VentureBeat rapportait en 2019 que 87\,\% des projets de data science
n'atteignent jamais la production. Ce n'est pas un probl\`eme de mod\'elisation
--- c'est un probl\`eme d'\textbf{ing\'enierie}.

\begin{center}
\begin{tikzpicture}[
  block/.style={rectangle, draw, rounded corners, minimum width=2.8cm,
    minimum height=1cm, align=center, font=\small},
  arrow/.style={-{Stealth[length=3mm]}, thick},
  fail/.style={-{Stealth[length=3mm]}, thick, red!70!black, dashed}
]
  \node[block, fill=blue!10] (data) {Donn\'ees};
  \node[block, fill=blue!10, right=1.2cm of data] (notebook) {Notebook\\Exp\'erimentation};
  \node[block, fill=green!10, right=1.2cm of notebook] (model) {Mod\`ele\\entra\^in\'e};
  \node[block, fill=orange!10, right=1.2cm of model] (deploy) {D\'eploiement\\Production};
  \node[block, fill=red!10, right=1.2cm of deploy] (monitor) {Monitoring\\Maintenance};

  \draw[arrow] (data) -- (notebook);
  \draw[arrow] (notebook) -- (model);
  \draw[fail] (model) -- node[above, font=\scriptsize, red!70!black] {\'Echec fr\'equent} (deploy);
  \draw[arrow] (deploy) -- (monitor);

  % Valley of death
  \node[below=0.6cm of model, font=\small\itshape, text=red!70!black] {<<~Vall\'ee de la mort~>>};
\end{tikzpicture}
\end{center}

Les causes principales de ces \'echecs sont~:
\begin{itemize}
  \item \textbf{Dette technique}\index{dette technique}~: le code ML accumule
    de la complexit\'e cach\'ee (d\'ependances de donn\'ees, boucles de
    r\'etroaction, configurations non versionn\'ees)~\cite{sculley2015}.
  \item \textbf{Irr\'eproductibilit\'e}~: impossible de recr\'eer les r\'esultats
    d'un coll\`egue, m\^eme avec le m\^eme code, car l'environnement, les donn\'ees
    ou les hyperparamm\`etres diff\`erent.
  \item \textbf{Foss\'e organisationnel}~: les data scientists travaillent en
    notebooks isol\'es, les ing\'enieurs ML ne comprennent pas les choix de
    mod\'elisation, et les \'equipes ops ne savent pas surveiller un mod\`ele.
  \item \textbf{Absence de monitoring}~: un mod\`ele d\'eploy\'e sans surveillance
    d\'erive silencieusement (\emph{data drift}, \emph{concept drift}).
  \item \textbf{Pas de pipeline reproductible}~: chaque \'etape (pr\'etraitement,
    entra\^inement, \'evaluation) est ex\'ecut\'ee manuellement.
\end{itemize}

\section*{Ce que le MLOps r\'esout}

Le \textbf{MLOps}\index{MLOps} (Machine Learning Operations) est un ensemble
de pratiques qui unifient le d\'eveloppement ML (Dev) et les op\'erations (Ops)
pour automatiser, standardiser et fiabiliser le cycle de vie des mod\`eles.
C'est l'\'equivalent du DevOps pour le machine learning, avec des d\'efis
sp\'ecifiques~: versionnement des donn\'ees, suivi des exp\'eriences,
reproductibilit\'e des r\'esultats, monitoring de la d\'erive.

\begin{center}
\begin{tikzpicture}[
  phase/.style={rectangle, draw, rounded corners, minimum width=2.2cm,
    minimum height=0.9cm, align=center, font=\small\bfseries},
  arrow/.style={-{Stealth[length=2.5mm]}, thick, steelblue}
]
  \node[phase, fill=blue!15] (design) {Conception};
  \node[phase, fill=blue!15, right=0.8cm of design] (dev) {D\'eveloppement};
  \node[phase, fill=green!15, right=0.8cm of dev] (train) {Entra\^inement};
  \node[phase, fill=orange!15, right=0.8cm of train] (deploy) {D\'eploiement};
  \node[phase, fill=red!15, right=0.8cm of deploy] (monitor) {Monitoring};

  \draw[arrow] (design) -- (dev);
  \draw[arrow] (dev) -- (train);
  \draw[arrow] (train) -- (deploy);
  \draw[arrow] (deploy) -- (monitor);
  \draw[arrow, bend right=30] (monitor) to node[below, font=\scriptsize] {r\'e-entra\^inement} (train);
\end{tikzpicture}
\end{center}

\section*{Organisation de ce cours}

Ce cours est con\c{c}u pour des \'etudiants de niveau Master en science des
donn\'ees, apprentissage automatique ou math\'ematiques appliqu\'ees. Il
suppose des bases solides en Python, une familiarit\'e avec la ligne de
commande Linux/Unix, et des connaissances en apprentissage automatique.

Chaque chapitre suit une structure p\'edagogique coh\'erente~:
\begin{enumerate}
  \item \textbf{Motivation et concepts}~: pourquoi cet outil ou cette pratique
    est n\'ecessaire.
  \item \textbf{Tutoriel pratique}~: code fonctionnel que l'\'etudiant peut
    ex\'ecuter sur sa machine.
  \item \textbf{Bonnes pratiques et anti-patterns}~: ce qu'il faut faire et
    ce qu'il faut \'eviter.
  \item \textbf{Comparaison d'outils}~: tableaux comparatifs quand plusieurs
    solutions existent.
  \item \textbf{Mini-projet}~: exercice int\'egrateur utilisant les outils des
    chapitres pr\'ec\'edents.
  \item \textbf{Exercices}~: gradu\'es par difficult\'e
    ($\bigstar$ configuration, $\bigstar\bigstar$ construction de pipeline,
    $\bigstar\bigstar\bigstar$ projet de bout en bout).
  \item \textbf{Aide-m\'emoire}~: r\'esum\'e des commandes essentielles.
\end{enumerate}

\medskip
\noindent\textbf{Remarque terminologique.}
Le MLOps emprunte massivement au vocabulaire anglophone. Quand un terme
technique n'a pas d'\'equivalent fran\c{c}ais \'etabli (par exemple
\emph{feature store}, \emph{drift}, \emph{pipeline}), nous l'utilisons en
anglais, en italique lors de sa premi\`ere apparition, avec une explication
en fran\c{c}ais.

\vfill
\noindent\textit{Bonne lecture et bons d\'eploiements~!}
